\documentclass[12pt, a4paper]{article}

%  Alguns pacotes iniciais
\usepackage[utf8]{inputenc}   % Codificação de caracteres (acentuação direta)
\usepackage[T1]{fontenc}      % Codificação de fontes para o português
\usepackage[brazil]{babel}    % Tradução de termos e hifenização para o português
\usepackage{geometry}         % Definição de margens
\geometry{top=3cm, bottom=2cm, left=3cm, right=2cm}

%  Pacotes de notação matemática
\usepackage{amsmath}   % Ferramentas avançadas para matemática
\usepackage{amsfonts}  % Conjuntos de números (reais, inteiros, etc.)
\usepackage{amssymb}   % Símbolos matemáticos adicionais

%  Cabeçalho
\title{Título}
\author{Seu Nome \\ 
\small Universidade Federal}
\date{\today}

\begin{document}

\maketitle % Aqui gera-se um título com base no cabeçalho

\begin{abstract}
Aqui é o resumo do artigo. O LaTeX separa a estrutura lógica do texto da formatação visual, permitindo que o pesquisador foque no conteúdo.
\end{abstract}

% Explicar aqui as seções
\section{Introdução}
Aqui começa a introdução do seu texto acadêmico. O LaTeX lida de forma automática com a numeração de seções.

% E subseções
\subsection{Justificativa}
Subseções.

\section{Notações de fórmulas  matemáticas}

Texto corrido pode conter equações inline, como $y_i = \beta_0 + \beta_1 x_i + \epsilon_i + \delta_i +\gamma_i$, que se ajustam perfeitamente à linha.

Para equações destacadas (display mode) e numeradas, utilizamos o ambiente \textbf{equation}:
\begin{equation}
    \hat{\beta}_i = \frac{\sum_{i=1}^{n} (x_i - \bar{x})(y_i - \bar{y})}{\sum_{i=1}^{n} (x_i - \bar{x})^2}
    \label{eq:mqo}
\end{equation}

Podemos referenciar a Equação \ref{eq:mqo} automaticamente ao longo do texto usando os comandos \texttt{\textbackslash label} e \texttt{\textbackslash ref}.

Para sistemas de equações ou alinhamento complexo, usamos o ambiente \textbf{align*}:
\begin{align*}
    f(x) &= (2x + 3)/2x \\
    f(x) &=\frac{2x+3}{2x} \\
    f(x) + g(x) &= x^2 + 2x + 2
\end{align*}

%Exemplo da função dentro do \[    \] no lugar do & e sem precisar usar o align
    \[g(x)= x^2 - 1\] 

Se eu não quiser alinhar usa-se \textbf{enumerate}: \\

\begin{enumerate}
%Exemplo da função dentro do $ $
\item[A)] $f(x)=\frac{2x+3}{2x}$ \\
\item[B)] $f(x)=x^2 - 1$ \\
\item [C)] $fxf(x) + g(x)= x^2 + 2x + 2$


\end{enumerate}

\end{document}